% subsubsection: 逆数型  \, $\displaystyle a_{n+1}=\frac{pa_n}{qa_n+r}$

\subsubsection{\texorpdfstring{逆数型   $\displaystyle a_{n+1}=\frac{pa_n}{qa_n+r}$}{逆数型}}\label{節：逆数型}
分子が$a_n$の定数倍である\bookterm{recurrence}{漸化式}を考えよう.
逆数を取れば,分母にあった一次式を$a_n$で割り,定数と$\frac1{a_n}$の和に分けられる.
以下では,$p\ne0$かつ$a_n\ne0$であり,$qa_n+r\ne0$として議論する.\footnote{$a_1=0$は逆数を取る前に分ける.$r\ne0$ならすべての項が$0$となる.$r=0$なら最初の更新から分母が$0$となり,そのままでは定義できない.}
両辺の逆数を取り,$b_n=\dfrac{1}{a_n}$と置くと,
\begin{align*}
  \frac{1}{a_{n+1}}&=\frac{qa_n+r}{pa_n}
  =\frac{q}{p}+\frac{r}{p}\cdot\frac{1}{a_n},\\
  b_{n+1}&=\frac{r}{p}b_n+\frac{q}{p}
\end{align*}
ここから特性方程式型または等差型として$b_n$を求め,最後に$a_n=\dfrac{1}{b_n}$へ戻す.
ただし$b_n=0$となる項は逆数を戻せないため,元の漸化式が定義されるかも確認する.

\noindent \bookblocklink{example-reciprocal-1}{problem}{answer}{【例題】}\par\nobreak
\begin{enumerate}[label=(\arabic*),parsep=3pt,format=\bookitemlink{example-reciprocal-1}{problem}{answer}]
  \item $\displaystyle a_1=1,\, a_{n+1}=\frac{a_n}{a_n+1}$
  \item $\displaystyle a_1=1,\, a_{n+1}=\frac{-a_n}{2a_n+1}$
\end{enumerate}

\noindent\bookblocklink{example-reciprocal-1}{hint}{problem}{【方針】}\par\nobreak
\begin{enumerate}[label=(\arabic*),parsep=3pt,format=\bookitemlink{example-reciprocal-1}{hint}{problem}]
  \item 逆数を取ると,$\dfrac{1}{a_n}$が毎回$1$ずつ増える形になります.
  \item 逆数を取って一次の\bookterm{recurrence}{漸化式}に帰着させましょう.
\end{enumerate}

\noindent\bookblocklink{example-reciprocal-1}{answer}{problem}{【解答】}\par\nobreak
\begin{enumerate}[label=(\arabic*),parsep=3pt,format=\bookitemlink{example-reciprocal-1}{answer}{problem}]
  \item $\displaystyle a_n=\frac{1}{n}$
  \item $\displaystyle a_n=\frac{1}{2(-1)^{n-1}-1}$
\end{enumerate}

\noindent\bookblocklink{example-reciprocal-1}{solution}{problem}{【解法】}\par\nobreak
\begin{enumerate}[label=(\arabic*),parsep=3pt,format=\bookitemlink{example-reciprocal-1}{solution}{problem}]
  \item $\displaystyle a_1=1,\, a_{n+1}=\frac{a_n}{a_n+1}$\\
    正の項と分母が続くので,逆数を取れる.$b_n=\dfrac{1}{a_n}$と置くと,
    \begin{align*}
      &\frac{1}{a_{n+1}}=1+\frac{1}{a_n},\\
      &b_{n+1}=b_n+1,\quad b_1=1\\
      &\therefore \quad b_n=n,\quad a_n=\frac1n
    \end{align*}
  \item $\displaystyle a_1=1,\, a_{n+1}=\frac{-a_n}{2a_n+1}$\\
    初めの更新は$1\longmapsto-\dfrac{1}{3}\longmapsto1$となる.
    項も分母も零でないので,$b_n=\dfrac{1}{a_n}$と置くと,
    \begin{align*}
      &\frac{1}{a_{n+1}}=-2-\frac{1}{a_n},\\
      &b_{n+1}=-b_n-2,\quad b_1=1
    \end{align*}
    \bookterm{fixed}{不動点}は$-1$なので,$c_n=b_n+1$と置き,定数項を消す.
    \begin{align*}
      &c_{n+1}=-c_n,\quad c_1=2\\
      &\therefore \quad b_n=c_n-1=2(-1)^{n-1}-1,\\
      &a_n=\frac{1}{2(-1)^{n-1}-1}
    \end{align*}
\end{enumerate}
